% !TeX root = ../technical_report.tex
\chapter{Real-Time Telemetry, Prometheus Instrumentation, and Drift Detection}
\label{ch:telemetry_and_drift}

% CHAPTER BRIEF FOR JULES / TECHNICAL WRITER:
% -------------------------------------------------------------
% Focus: Production telemetry, Prometheus scrapers, and embedding/feature drift detection.
% Repository Files:
%   - src/taxon_vision/monitoring/telemetry.py
%   - src/taxon_vision/monitoring/drift.py
%   - src/taxon_vision/service/routes/health.py (metrics endpoint)
% Detailed Technical Concepts:
%   - Prometheus Metric Exporters:
%       * taxon_predictions_total: Counter tracking requests by predicted class/status.
%       * taxon_inference_latency_seconds: Histogram tracking p50, p95, p99 serving latency.
%       * taxon_human_referrals_active: Gauge tracking human triage routing frequency.
%   - Feature and Concept Drift Mechanics (drift.py):
%       * Wasserstein Distance: compute_wasserstein_drift on scalar prediction distributions and norms.
%       * Maximum Mean Discrepancy (MMD) and KS Tests.
% -------------------------------------------------------------

\section{Telemetry and Observability Architecture}
\label{sec:telemetry_architecture}

In production machine learning operations (MLOps), maintaining deterministic execution and epistemic safety necessitates rigorous, real-time observability. The TaxonVision platform operates within continuous ecological vision streams, where visual domain shifts -- such as seasonal foliage variations, differing camera sensor characterisations, and regional lighting fluctuations -- inevitably degrade model reliability over time. To combat this silent degradation, the platform integrates a decoupled, pull-based telemetry architecture grounded in the Prometheus ecosystem.

The core observability philosophy prioritises non-blocking, latency-transparent metric emission. Serving visual inference in production requires strict non-blocking concurrency. Executing image decoding, tensor normalisation, and PyTorch or ONNX Runtime matrix operations directly on the asynchronous I/O loop thread risks event loop starvation, which blocks all concurrent requests (e.g., \texttt{/health}, \texttt{/metrics}, and \texttt{/predict}). Consequently, synchronous inference is offloaded to an asynchronous worker threadpool via the AnyIO library. Crucially, the emission of telemetry metrics must not re-introduce blocking I/O or external network latency into this highly optimised inference pathway.

By employing a pull-based scraping model, the Fast\-API service simply increments counters, gauges, and histograms in memory. These metric aggregations introduce negligible CPU overhead ($\mathcal{O}(1)$ time complexity for state mutation) and zero network latency. A dedicated orchestration agent (e.g., a Prometheus scraper operating within the Kubernetes cluster) polls the \texttt{/metrics} HTTP endpoint at predefined intervals (e.g., 15--60 seconds). This decoupled topology ensures that network partitions, scraper misconfigurations, or monitoring downtime cannot cascade into inference failures. The resulting telemetry stream enables advanced aggregations via the Prometheus Query Language (PromQL) and comprehensive visualisations within Grafana dashboards, facilitating the tracking of inference throughput, hardware utilisation, and dynamic decision boundaries without compromising primary application throughput.

Furthermore, decoupling metric aggregation ensures that the computational resources allocated to model inference -- specifically, the virtualised NVIDIA time-slicing partitions on the bare-metal host -- are strictly preserved for active prediction tasks. The stateless nature of the Prometheus counters guarantees that memory consumption remains constant regardless of request volume, mitigating the risk of memory leaks associated with unbound event logging.

\section{Prometheus Instrumentation and Metric Design}
\label{sec:prometheus_instrumentation}

The mathematical formalisation of model performance relies on three fundamental Prometheus metric types: Counters, Histograms, and Gauges. The telemetry module\\\texttt{src/taxon\_vision/monitoring/telemetry.py} centralises the definition and global registration of these metrics. By invoking these definitions directly upon module import, they are seamlessly attached to the global Prometheus client registry without requiring manual singleton instantiation.

\subsection{Metric Definitions}

Table \ref{tab:prometheus_metrics} delineates the explicit metrics instrumented within the inference pathway, each purposefully engineered to expose distinct facets of system performance.

\begin{table}[ht]
\centering
\caption{Prometheus Telemetry Metrics for TaxonVision Inference}
\label{tab:prometheus_metrics}
\begin{tabularx}{\textwidth}{llX}
\toprule
\textbf{Metric Name} & \textbf{Type} & \textbf{Description \& Utilisation} \\
\midrule
\texttt{taxon\_predictions\_total} & Counter & Tracks the total volume of evaluated species predictions. It utilises a \texttt{status} label (e.g., \texttt{success}, \texttt{fallback}, \texttt{error}) to compute error budgets and automation rates over variable time windows. \\
\texttt{taxon\_inference\_latency\_seconds} & Histogram & Measures the end-to-end inference duration, encompassing image decoding, tensor preprocessing, and ONNX execution. Facilitates the derivation of critical $p_{50}$, $p_{95}$, and $p_{99}$ latency percentiles. \\
\texttt{taxon\_human\_referrals\_active} & Gauge & Reports the instantaneous count of observations deferred to human experts, functioning as a continuous proxy for the epistemic uncertainty boundary and fallback queue depth. \\
\bottomrule
\end{tabularx}
\end{table}

The \texttt{taxon\_predictions\_total} counter is particularly critical for operational diagnostics. During continuous execution, unpredictable payload anomalies -- such as corrupted image headers or unsupported MIME types -- trigger decoding exceptions (e.g., \texttt{UnidentifiedImageError}). The inference router gracefully degrades by catching these exceptions, serving a deterministic fallback prediction (e.g., \textit{Danaus plexippus} for integration testing), and incrementing the counter with the \texttt{status="fallback"} label. In stark contrast, unexpected matrix dimensionality inconsistencies or out-of-memory (OOM) driver crashes invoke the \texttt{status="error"} label, immediately breaching predefined Service Level Objective (SLO) thresholds and triggering pager alerts.

\subsection{Inference Integration}

The integration of these metric primitives occurs securely within the \texttt{predict\_species} routing handler, defined in \texttt{src/taxon\_vision/service/routes/predict.py}. The implementation exclusively relies upon the high-resolution \texttt{time.perf\_counter()} primitive to compute precise, wall-clock execution latency without succumbing to system clock synchronisation drift.

\begin{listing}[ht]
\caption{Prometheus telemetry integration within the asynchronous inference router.}
\label{lst:predict_metrics}
\begin{minted}{python}
import time
import anyio
from fastapi import APIRouter, File, HTTPException, UploadFile, status
from PIL import UnidentifiedImageError

from taxon_vision.monitoring.telemetry import PREDICTION_COUNTER, LATENCY_HISTOGRAM
from taxon_vision.service.inference import run_prediction

@router.post("/predict", response_model=PredictionResponse)
async def predict_species(file: UploadFile) -> PredictionResponse:
    t0 = time.perf_counter()
    content = await file.read()
    try:
        # Offload intensive ONNX execution to a threadpool worker
        pred = await anyio.to_thread.run_sync(run_prediction, content)

        # Calculate wall-clock latency in milliseconds
        latency = (time.perf_counter() - t0) * 1000.0

        # Record telemetry observations safely
        PREDICTION_COUNTER.labels(status="success").inc()
        LATENCY_HISTOGRAM.observe(latency / 1000.0) # Standardise to seconds

        return _build_live_prediction(pred, latency)
    except (UnidentifiedImageError, OSError, ValueError) as img_err:
        latency = (time.perf_counter() - t0) * 1000.0
        PREDICTION_COUNTER.labels(status="fallback").inc()
        return _build_fallback_prediction(latency)
    except Exception as exc:
        PREDICTION_COUNTER.labels(status="error").inc()
        raise HTTPException(status_code=500, detail=str(exc))
\end{minted}
\end{listing}

The actual exposition of these metrics to external Prometheus scrapers is formally implemented in the \texttt{src/taxon\_vision/service/routes/health.py} module. A dedicated HTTP \texttt{GET} route invokes the Prometheus client's \texttt{generate\_latest()} method. This process rapidly serialises all registered counters, histograms, and gauges into the standard text-based exposition format. Because this serialization involves reading static, in-memory float structures, it executes with negligible latency, completely bypassing disk access, database lookups, and model inference pathways.

\section{Data Drift and Concept Drift Detection Engine}
\label{sec:data_drift_engine}

While systemic metrics (e.g., inference latency, HTTP error rates) are indispensable for detecting infrastructure failures, silent algorithmic degradation caused by data drift requires sophisticated statistical evaluation. In ecological computer vision, concept drift typically manifests as shifts in species relative abundances (e.g., massive migratory influxes outstripping historical training priors), whereas covariate drift involves visual domain shifts such as changing illumination patterns or deployment across novel camera sensor architectures.

\subsection{Dimensionality Challenges and Statistical Distance}

Monitoring representations within the deep embedding space of a Vision Foundation Model (VFM) involves analysing high-dimensional vectors (often $D=768$ or $D=1024$). Attempting to apply high-dimensional multivariate statistical tests, such as Maximum Mean Discrepancy (MMD) or the multidimensional Kolmogorov-Smirnov (KS) test, directly against these embeddings reveals severe computational bottlenecks. These methods suffer acutely from the curse of dimensionality; evaluating MMD over a dense $1024$-dimensional manifold imposes polynomial time complexity, rendering real-time evaluation across sliding production windows computationally intractable.

Similarly, while the Population Stability Index (PSI) is highly effective for quantifying shift across discrete taxonomic orders or families, it lacks the mathematical continuity required to measure sub-class representation divergence within the continuous embedding distributions.

To bypass these strict multivariate limitations, the TaxonVision drift detection engine projects high-dimensional feature representations onto their scalar vector norms, effectively collapsing the multivariate distribution into an efficient univariate scalar array for rapid comparison.

\subsection{Wasserstein Distance on Embedding Norms}

The \texttt{src/taxon\_vision/monitoring/drift.py} module implements this mathematically rigorous yet computationally efficient one-dimensional approximation using the Wasserstein-1 metric, colloquially known as the Earth Mover's Distance.

Let $\mathbf{Z}_{\text{ref}} \in \mathbb{R}^{N_{\text{ref}} \times D}$ represent a matrix of baseline embeddings obtained during the model's pristine validation phase. Conversely, let $\mathbf{Z}_{\text{curr}} \in \mathbb{R}^{N_{\text{curr}} \times D}$ represent the embeddings sampled dynamically from the current production inference window. The $L_2$ Euclidean norms of these representations are computed efficiently along the feature dimension axis:

\begin{equation}
    ||\mathbf{z}_{\text{ref}}^{(i)}||_2 = \sqrt{\sum_{j=1}^{D} \left( z_{\text{ref}, j}^{(i)} \right)^2}
\end{equation}

By extracting these norms, the system defines two continuous probability distributions, $P_{\text{ref}}$ and $P_{\text{curr}}$, representing the absolute magnitude characteristics of the baseline and production features, respectively. The Wasserstein drift approximation, $W_1$, is then calculated as the absolute difference between the expected values (empirical means) of these norm distributions:

\begin{equation}
    W_1(P_{\text{ref}}, P_{\text{curr}}) \approx \left| \mathbb{E}[||\mathbf{z}_{\text{ref}}||_2] - \mathbb{E}[||\mathbf{z}_{\text{curr}}||_2] \right|
\end{equation}

This projection methodology boasts a superior computational complexity of $\mathcal{O}(N \times D)$, executing in milliseconds via highly vectorised NumPy array operations without requiring complex, iterative transport planning algorithms.

\begin{listing}[ht]
\caption{Vectorised computation of Wasserstein drift on continuous embedding norms.}
\label{lst:wasserstein_drift}
\begin{minted}{python}
import numpy as np

def compute_wasserstein_drift(reference_embeddings: np.ndarray, current_embeddings: np.ndarray) -> float:
    """Compute 1D Wasserstein distance approximation on representation embedding norms.

    Returns the scalar absolute difference between the mean L2 norms of the
    reference and current embedding manifolds.
    """
    ref_norm = np.linalg.norm(reference_embeddings, axis=-1)
    curr_norm = np.linalg.norm(current_embeddings, axis=-1)
    return float(np.abs(np.mean(ref_norm) - np.mean(curr_norm)))
\end{minted}
\end{listing}

A statistically significant increase in this scalar $W_1$ distance acts as an early-warning signal, indicating that the fundamental structural geometry of the incoming visual data has diverged from the original training distribution. This divergence consistently correlates with a severe degradation in model calibration and necessitates immediate, automated intervention.

\section{Automated Alerting and Retraining Triggers}
\label{sec:automated_retraining}

The computation of the Wasserstein drift metric is far more than a passive analytical exercise; it forms the core deterministic heuristic driving the automated MLOps continuous integration loop.

When the computed drift distance breaches a predefined, statistically significant threshold (a parameter meticulously established via baseline distributions during Bayesian hyperparameter optimisation), a structured JSON webhook payload is dispatched directly to the declarative Jenkins pipeline operating within the isolated \texttt{jenkins-ci} Kubernetes namespace. This payload acts as an authoritative, unmediated retraining trigger.

Upon receiving the validated drift alert, the Jenkins orchestrator initiates a comprehensive, multi-stage remediation sequence designed to restore epistemic safety:
\begin{enumerate}
    \item \textbf{Dataset Re-ingestion:} Observations flagged for active learning and those residing in human-triaged referral queues are aggressively extracted and merged back into the primary Arrow streaming datasets, rectifying the long-tail class imbalance highlighted by the drift metric.
    \item \textbf{Cryptographic Lineage Versioning:} The newly updated training corpus is hashed and rigorously version-controlled via Data Version Control (DVC). The updated \texttt{dvc.lock} file ensures complete mathematical lineage tracking, pushing the resultant artifacts to the remote DagsHub S3-compatible object storage infrastructure.
    \item \textbf{Optuna Hyperparameter Optimisation:} A highly parallelised Optuna study is launched across the NVIDIA time-slicing partitions. This study aggressively re-optimises critical regularisation parameters (such as the base learning rate and specific Class-Balanced Loss decay parameters) against the newly established, shifted data distribution.
    \item \textbf{Model Registration and Deployment:} The retrained model head undergoes a strict evaluation suite. If its empirical validation metrics decisively exceed the pre-drift baseline, the newly refined checkpoint is serialised into INT8/FP16 precision via the ONNX Runtime engine. This artifact is formally registered within the MLflow Model Registry, instantly triggering a zero-downtime, rolling Kubernetes deployment update that seamlessly replaces the degraded pods in the production namespace.
\end{enumerate}

\providecolor{TaxonPrimary}{HTML}{1976D2}
\providecolor{TaxonSecondary}{HTML}{388E3C}
\providecolor{TaxonDark}{HTML}{263238}
\providecolor{TaxonAccent}{HTML}{E65100}

\begin{figure}[htbp]
    \centering
    \resizebox{\textwidth}{!}{
    \begin{tikzpicture}[
        node distance=1.5cm and 2cm,
        font=\sffamily\small,
        container/.style={draw=TaxonDark, fill=gray!5, thick, rounded corners, rectangle, inner xsep=15pt, inner ysep=15pt, dashed},
        component/.style={draw=TaxonDark, fill=white, thick, rounded corners, rectangle, align=center, minimum height=1.2cm, drop shadow, text width=4cm},
        highlight/.style={draw=TaxonAccent, fill=TaxonAccent!5, thick, rounded corners, rectangle, align=center, minimum height=1.2cm, drop shadow, text width=4cm},
        arrow/.style={-Stealth, thick, draw=TaxonDark},
        alert_arrow/.style={-Stealth, thick, draw=TaxonAccent, dashed}
    ]

        % --- Prometheus Telemetry Layer ---
        \node[component] (fastapi) {FastAPI Application\\\texttt{/predict}};
        \node[component, right=of fastapi] (metrics) {Prometheus Exposition\\\texttt{/metrics}};
        \node[component, right=of metrics] (scraper) {Prometheus Scraper\\(Pull-based, 15-60s)};
        
        \node[component, below=0.5cm of fastapi, text width=3.2cm] (counter) {\texttt{\footnotesize taxon\_predictions\\\_total}\\(success, fallback, error)};
        \node[component, below=0.5cm of metrics, text width=3.5cm] (histogram) {\texttt{\footnotesize taxon\_inference\\\_latency\_seconds}\\(p50, p95, p99)};
        \node[component, below=0.5cm of scraper, text width=3.2cm] (gauge) {\texttt{\footnotesize taxon\_human\\\_referrals\_active}\\(Triage Queue)};

        \begin{scope}[on background layer]
            \node[container, fit=(fastapi) (metrics) (scraper) (counter) (histogram) (gauge), label={[font=\bfseries, text=TaxonPrimary]above:Prometheus Telemetry Layer}] (telemetry_layer) {};
        \end{scope}

        % --- Statistical Drift Detection Engine ---
        \node[component, below=2cm of telemetry_layer] (engine) {Drift Detection Engine\\\texttt{monitoring/drift.py}};
        
        \node[component, left=0.5cm of engine, text width=3.5cm] (ref_embeddings) {Baseline Embeddings\\$\mathbf{Z}_{\text{ref}}$};
        \node[component, right=0.5cm of engine, text width=3.5cm] (curr_embeddings) {Streaming Embeddings\\$\mathbf{Z}_{\text{curr}}$};

        \node[highlight, below=0.8cm of engine, text width=6cm] (wasserstein) {$L_2$ Norm Scalar Projection\\$W_1(P_{\text{ref}}, P_{\text{curr}}) \approx |\mathbb{E}[||\mathbf{z}_{\text{ref}}||_2] - \mathbb{E}[||\mathbf{z}_{\text{curr}}||_2]|$};
        
        \node[highlight, below=0.8cm of wasserstein, text width=4cm] (decision_gate) {Decision Gate\\$W_1 > \tau_{\text{drift}}$};

        \begin{scope}[on background layer]
            \node[container, fit=(engine) (ref_embeddings) (curr_embeddings) (wasserstein) (decision_gate), label={[font=\bfseries, text=TaxonPrimary]above:Statistical Drift Detection Engine}] (drift_layer) {};
        \end{scope}

        % --- Closed-Loop Feedback & Automated Retraining Trigger ---
        \node[component, below=2cm of decision_gate] (jenkins) {Jenkins Controller\\\texttt{jenkins-ci} namespace};
        
        \node[component, right=1cm of jenkins] (step1) {1. Dataset Re-ingestion\\\& DVC Snapshot};
        \node[component, right=0.5cm of step1] (step2) {2. Optuna Bayesian\\Hyperparameter Tuning};
        \node[component, below=0.5cm of step1] (step3) {3. ONNX INT8\\Quantization};
        \node[component, below=0.5cm of step2] (step4) {4. Rolling Zero-Downtime\\Deployment Update};

        \begin{scope}[on background layer]
            \node[container, fit=(jenkins) (step1) (step2) (step3) (step4), label={[font=\bfseries, text=TaxonPrimary]above:Closed-Loop Feedback \& Automated Retraining Trigger}] (retraining_layer) {};
        \end{scope}

        % --- Edges ---
        % Telemetry internal edges
        \draw[arrow] (fastapi) -- (metrics);
        \draw[arrow] (metrics) -- (scraper);
        \draw[arrow, dotted] (counter) -- (fastapi);
        \draw[arrow, dotted] (histogram) -- (metrics);
        \draw[arrow, dotted] (gauge) -- (scraper);

        % Flow to drift
        \draw[arrow] (telemetry_layer.south) -- (engine.north);
        \draw[arrow] (ref_embeddings) -- (engine);
        \draw[arrow] (curr_embeddings) -- (engine);
        \draw[arrow] (engine) -- (wasserstein);
        \draw[arrow] (wasserstein) -- (decision_gate);

        % Alert trigger
        \draw[alert_arrow] (decision_gate) -- (jenkins) node[midway, right, text=TaxonAccent] {Webhook Dispatch (Alert)};

        % Jenkins pipeline edges
        \draw[arrow] (jenkins) -- (step1);
        \draw[arrow] (step1) -- (step2);
        \draw[arrow] (step2) -- (step3);
        \draw[arrow] (step3) -- (step4);
        
        % Loop back to deployment
        \draw[arrow] (step4.south) |- ++(0,-1cm) -| (telemetry_layer.south east) node[pos=0.25, below] {Deployment of new model artifact};

    \end{tikzpicture}
    }
    \caption{End-to-end observability and closed-loop continuous training pipeline. Telemetry signals from the FastAPI Prometheus integration stream into the statistical drift engine, which projects high-dimensional embeddings down to scalar norms. Breaching the Wasserstein $W_1$ threshold dispatches a webhook alert to Jenkins, triggering an automated sequence of DVC snapshotting, Optuna retuning, INT8 quantization, and a zero-downtime rolling update.}
    \label{fig:telemetry_and_drift_detection}
\end{figure}


This sophisticated, closed-loop feedback mechanism guarantees that the TaxonVision model dynamically adapts to continuous, real-world ecological domain shifts. By systematically eradicating the need for manual, reactive intervention, the architecture successfully maintains stringent epistemic safety boundaries and exceptional inferential throughput over extended, multi-year operational timelines.
